\section{数值研究}
\label{sec:numerical}




为在数值上检验上述想法，我们在 $\beta=6.0$、体积 $32^3\times 64 $
的格点（格距 $a=0.093$ fm）上进行了淬火近似下的格点 QCD 计算。
费米子作用量采用非微扰调好的 clover 作用量，
clover 系数由 Alpha 合作组计算 \cite{Luscher:1996ug}。

我们共使用 500 个组态，相邻组态间隔 1000 次更新，
每次更新包含四次过弛豫（over-relaxation）扫描和一次热浴（heatbath）扫描。
每个组态上用 6 个随机选取的点源计算关联函数。
该设定下 π 介子与核子质量分别测定为 $601(1)$ MeV
与 $1411(4)$ MeV。到物理能量单位的换算采用了 Alpha 合作组的
淬火 QCD 标度设定 \cite{Necco:2001xg}。

我们的核子态被助推到总动量最高 $2.5\, $GeV
（相应于第 6 个格点动量）。
在此动量范围内，核子的连续色散关系在计算误差范围内得到满足，
表明 ${\cal O}(aP)$ 的格点伪影很小。
图 \ref{fig:disp} 画出了核子能量随动量的变化，
并连同与我们格点零动量核子能量相对应的连续色散关系。

    \begin{figure}[t]
  \includegraphics[width=0.55\textwidth]{images/disperkey.pdf}
  \caption{核子色散关系。能量与动量均以格点单位表示。实线是连续色散关系（并非拟合），误差带指示格点核子能量的统计误差。}
  \label{fig:disp}
\end{figure}

矩阵元的计算采用文献~\cite{Bouchard:2016heu} 所述的方法，
算符插入由式~(\ref{Malpha}) 给出。
取流的分量，我们便可分离出 $\mathcal{M}_p(-z\cdot p, -z^2)$，
如前所述它与 PDF 直接相关。

按照文献~\cite{Bouchard:2016heu}，需要计算两类关联函数。
第一类是通常的核子两点函数
%
 \begin{equation}
 C_P(t) = \langle \mathcal{N}_P(t)\overline{ \mathcal{N}}_P(0) \rangle\,  ,
 \end{equation}
%
其中 $ \mathcal{N}_P(t)$ 是具有动量 $p$ 的、螺旋度平均的非相对论性核子内插场。
$ \mathcal{N}_p(t)$ 中的夸克场经过规范不变的 Gauss 涂抹（smearing）。
已知这种内插场与核子基态耦合良好（讨论见文献~\cite{Bouchard:2016heu}）。
夸克涂抹宽度经过优化，以便在我们计算的动量范围内与核子基态有良好交叠。
第二类关联函数为
 \begin{equation}
 C^{\mathcal{O}^0(z)}_P(t) =\sum_\tau \langle \mathcal{N}_P(t) \mathcal{O}^0(z,\tau)\overline{ \mathcal{N}}_P(0) \rangle\, ,
 \end{equation}
其中
  \begin{equation}
 \mathcal{O}^0(z,t)= \overline\psi(0,t) \gamma^0 \tau_3 \hat{E}(0,z;A)\psi(z,t)\,,
 \end{equation}
$\tau_3$ 是味 Pauli 矩阵。
质子动量与夸克场的位移都沿 $\hat z$ 轴方向（$\vec z = z_3 \hat z$ 与 $\vec p = P \hat z$）。
我们定义有效矩阵元
 \begin{equation}
 \mathcal{M}_{\rm eff}(z_3 P, z_3^2;t) = \frac{C^{\mathcal{O}^0(z)}_P(t+1)}{C_P(t+1)} - \frac{C^{\mathcal{O}^0(z)}_P(t)}{C_P(t) }  \  .
 \end{equation}
如文献~\cite{Bouchard:2016heu} 所示，矩阵元 $\mathcal{J}$ 可在大欧氏时间间隔处抽取为
 \begin{equation}
\frac{\mathcal{J}(z_3P, z_3^2)}{2 E}= \lim_{t\rightarrow\infty}  \mathcal{M}_{\rm eff}(z_3 P, z_3^2;t) \  ,
 \end{equation}
其中 $E$ 是核子的能量。
      \begin{figure}[t]
  \includegraphics[width=0.48\textwidth]{images/plateau-p2z4.pdf} \hfill \includegraphics[width=0.48\textwidth]{images/plateau-p3z8.pdf}
  \caption{用于抽取约化矩阵元的典型拟合。上图对应 $p=2\pi/L \cdot 2$、$z=4$，
  下图对应 $p=2\pi/L \cdot 3$、$z=8$，动量与坐标均以格点单位表示。
  }
  \label{fig:plateau}
\end{figure}
与传统顺序源方法不同，这种抽取矩阵元的方法允许利用核子产生与湮没算符的全部源—阱间距来计算矩阵元。

所得有效矩阵元带有激发态污染，按 $e^{-t \Delta E }$ 标度，
其中 $t$ 是核子产生与湮没算符的欧氏时间间隔，
$\Delta E$ 是核子第一激发态的质量间隙。此外，该方法还能以不增加计算量为代价，
计算对应于不同核子动量、自旋极化及味结构的全部核子矩阵元。

因此该方法的总计算成本低于顺序源方法的等价成本，
尤其是因为我们在方法上侧重于获得尽可能多的核子动量态。
这一方法最近已被成功用于单核子及多核子矩阵元计算~\cite{Berkowitz:2017gql,Tiburzi:2017iux,Shanahan:2017bgi}。

关于格点矩阵元的归一化，注意到当 $z_3=0$ 时矩阵元
 $\mathcal{M}(z_3 P, z_3^2)$ 对应局域矢量（同位旋矢量）流，因而应等于 1；
但在格点上由于格点伪影，事实并非如此。
为此引入重正化常数
 \begin{equation}
 Z_P = \frac{1}{ \left.\mathcal{J}(z_3 P, z_3^2)\right|_{z_3=0}}\, .
 \label{eq:renorm}
 \end{equation}
因子 $Z_P$ 应与 $P$ 无关。然而，同样由于格点伪影或潜在的拟合系统误差，事实并非如此。
因此我们对每个动量分别用各自的 $Z_P$ 因子重正化矩阵元，
从而充分利用最大统计相关性{以减小统计误差}，
并让格点伪影在比值中相消。
于是矩阵元通过如下比值抽取：
  \begin{equation}
\mathcal{M}(z_3 P, z_3^2)= \lim_{t\rightarrow\infty}  \frac{\mathcal{M}_{\rm eff}(z_3 P, z_3^2;t)}{\left.\mathcal{M}_{\rm eff}(z_3 P, z_3^2;t)\right|_{z_3=0}}  \  .
 \end{equation}
为了确定约化矩阵元 $\mathfrak{M}(\nu,z_3^2)$，我们引入双比值
 \begin{align}
\mathfrak{M}(\nu, z_3^2)=&  \lim_{t\rightarrow\infty}
 \frac{\mathcal{M}_{\rm eff}(z_3 P, z_3^2;t)}{\left.\mathcal{M}_{\rm eff}(z_3 P, z_3^2;t)\right|_{z_3=0}} \nn
 &\times
\frac{\left.\mathcal{M}_{\rm eff}(z_3  P, z_3^2;t)\right|_{P=0,z_3=0}}
{\left.\mathcal{M}_{\rm eff}(z_3  P, z_3^2;t)\right|_{P=0}}  \ ,
\label{eq:redMatElem}
 \end{align}
它按照式~(\ref{eq:renorm}) 完成矢量流的重正化。
实践中，无穷大 $t$ 极限通过在适当拟合区间内对常数作拟合获得。
在我们研究的所有情形中，平均每自由度 $\chi^2$ 为 ${\cal O}(1)$ 量级。
用于抽取约化矩阵元的典型拟合见图~\ref{fig:plateau}。
所有拟合均使用完整协方差矩阵进行，误差棒用 jackknife 方法确定。

这里指出，式~(\ref{eq:redMatElem}) 定义的约化矩阵元具有良好定义的连续极限，
不需要额外的重正化。该连续极限在固定 $\nu$ 与 $z^2$ 以及固定夸克质量下取得。

本计算沿 $z$ 轴使用了直到 $6\cdot 2\pi/L $ 的动量，
对应的物理动量约为 2.5\,GeV。
